% Block 1 of the slides (01_llm.tex) and the attention backup slide.
\section{How a large language model works}\label{sec:llm}

\onslides{tokens; next-token prediction; generation as a loop; pre-training as a
fit; post-training (instruction tuning, human feedback, verifiable rewards); two tables
of properties and their consequences; knowledge cutoff and the Julia bias.}

\subsection{Tokens}

A model never sees characters. A \emph{tokenizer} splits text into pieces from a fixed
vocabulary and maps each piece to an integer. Modern tokenizers learn the vocabulary
from data with byte-pair encoding \rf{sennrich}: start from single bytes and repeatedly
merge the most frequent adjacent pair into a new symbol, until the vocabulary has the
target size. Frequent words become one token; rare words, numbers, code and \LaTeX{}
split into several. Llama~3 uses $128$K tokens and reaches $3.94$ characters per token
on English text \rf{llama3}. An English word averages about five letters plus a space,
so one token is roughly $0.7$ words, consistent with the ``$0.75$ words per token'' on the
slide. The token numbers on the slide are illustrative.

Two consequences come up in questions. Letter-level tasks (counting the r's in a word,
reversing a string) are hard because the model does not see letters. Arithmetic on long
numbers is unreliable because digits are grouped into tokens in irregular ways; this is
one reason agents call a calculator or write a script.

\subsection{Next-token prediction}

Write the tokens so far as $x_1, \ldots, x_n$; this sequence is the \emph{context}. The
model is a function that returns a probability distribution over the vocabulary for the
next token,
\begin{equation}
  p_\theta(x_{n+1} \mid x_1, \ldots, x_n),
\end{equation}
where $\theta$ denotes the parameters of a neural network, a \emph{transformer}
\rf{vaswani}. Large models have $10^{11}$ to $10^{12}$ parameters; the largest Llama~3
model has $4.05 \times 10^{11}$ \rf{llama3}. The network outputs one real number per
vocabulary entry, a \emph{logit} $z_i$, and the distribution is the softmax
\begin{equation}
  p_\theta(x_{n+1} = i \mid x_1, \ldots, x_n)
    = \frac{\ee^{z_i/T_{\mathrm{s}}}}{\sum_j \ee^{z_j/T_{\mathrm{s}}}},
\end{equation}
with a \emph{sampling temperature} $T_{\mathrm{s}}$ set by the harness (the original
distribution is $T_{\mathrm{s}} = 1$). A physicist reads this as a Boltzmann
distribution with energies $-z_i$. Small $T_{\mathrm{s}}$ concentrates the probability
on the most likely token; large $T_{\mathrm{s}}$ flattens it.

\textbf{Generation is a loop.} Sample a token, append it to the context, compute the
distribution again. A long answer is thousands of passes. Two facts follow and are on the
slide: in each pass the context is the model's only input, and a token once written is
never revised. A model can correct itself only by writing more text, for example
``wait, that is wrong'', which is what reasoning models learn to do (below).

\textbf{Attention} (backup slide). Inside the transformer, each token is a vector. In
each layer, each token's vector is mapped linearly to a \emph{query}, a \emph{key} and a
\emph{value}; stacking these for all tokens gives matrices $Q$, $K$, $V$, and the layer
computes
\begin{equation}
  \mathrm{Attention}(Q, K, V) = \mathrm{softmax}\!\left(\frac{Q K^{\mathsf{T}}}{\sqrt{d}}\right) V,
\end{equation}
where $d$ is the dimension of the keys and the softmax acts row by row \rf{vaswani}. Row
$m$ of the result is a weighted average of the values of all tokens, weighted by how well
their keys match the query of token $m$. In a language model a mask restricts the
average to earlier tokens. The matrix $Q K^{\mathsf{T}}$ has one entry per pair of
tokens, so the cost grows with the square of the context length; this is one reason the
context window is finite. 3Blue1Brown's chapters 5 and 6 \rf{3b1b} are the best visual
explanation.

\subsection{Pre-training: a very large fit}

Pre-training minimises the cross-entropy of the next token over a corpus,
\begin{equation}
  \mathcal{L}(\theta) = -\sum_{n} \ln p_\theta(x_{n+1} \mid x_1, \ldots, x_n),
\end{equation}
where the sum runs over every position $n$ of every document. This is a maximum-likelihood
fit: the same principle as a least-squares fit with Gaussian errors, with $10^{12}$
parameters and $10^{13}$ data points. The largest Llama~3 model was pre-trained on
$15.6 \times 10^{12}$ tokens \rf{llama3}. The loss falls as a power law in model size,
data and compute \rf{kaplan}, and for a fixed compute budget model size and data should
grow together \rf{chinchilla}; these \emph{scaling laws} are why models became large.

What the fit learns is the distribution of text. Text about physics is mostly correct, so
the model learns a great deal of physics; but the objective rewards the \emph{form} of
plausible text and never checks its truth. That is the root of the first row of the
property table: a fluent wrong derivation and a fluent right one are both high-likelihood
text.

\textbf{Knowledge cutoff.} The corpus ends at some date. Packages released or changed
after it are unknown or known in an old version. The corpus also reflects what people
write: there is far more Python than Julia. The slide's example shows the typical
result. Asked for QuantumOptics.jl, the model writes \cmd{mesolve}, the name of the
master-equation solver in QuTiP (Python); the Julia package calls it
\cmd{timeevolution.master} \rf{qojulia,qojulia-docs}. The code looks right and fails at
the first run, which an agent notices and a chatbot user may not.

\subsection{Post-training: from text model to assistant}

A pre-trained model continues text; asked a question, it may continue with more
questions. Post-training turns it into an assistant in three stages.

\textbf{Instruction tuning} (supervised fine-tuning). Continue training on examples of
prompts with good answers written by people. The model imitates the format and tone.

\textbf{Reinforcement learning from human feedback} (RLHF). People compare pairs of
answers to the same prompt. A second network, the \emph{reward model}, is trained to
predict which answer people prefer, giving a scalar score $r(x, y)$ for prompt $x$ and
answer $y$. The language model, now viewed as a policy $\pi_\theta(y \mid x)$, is then
trained by reinforcement learning to maximise
\begin{equation}
  \mathbb{E}_{y \sim \pi_\theta}\left[r(x, y)\right]
    - \beta\, D_{\mathrm{KL}}\!\left(\pi_\theta \,\|\, \pi_{\mathrm{ref}}\right),
\end{equation}
where $\pi_{\mathrm{ref}}$ is the model before this stage, $D_{\mathrm{KL}}$ the
Kullback--Leibler divergence and $\beta > 0$ a weight that keeps the model from drifting
far from fluent text. The method comes from Christiano et al.\ \rf{christiano} and was
applied to language models in InstructGPT \rf{instructgpt}, where a 1.3-billion-parameter
model trained this way was preferred by human raters over the 175-billion-parameter
GPT-3. Variants replace human raters with a model that judges answers against written
principles \rf{constitutional}.

\textbf{Reinforcement learning with verifiable rewards} (RLVR). Replace the learned reward
model with a check: the final answer to a maths problem matches the known one, or the
generated code passes the tests. The term RLVR was introduced with Tülu~3
\rf{tulu3}; DeepSeek-R1 \rf{deepseekr1} showed that this kind of training alone makes a
model produce long chains of intermediate reasoning, check its steps and backtrack,
before giving an answer. These are the ``thinking'' tokens of current models. Earlier,
the same behaviour had been elicited by prompting (``chain of thought'') \rf{wei-cot}.

The key box on the slide summarises the incentives: the model was rewarded for answers
that \emph{look good to people} and \emph{pass checks}. Neither is the same as being
correct, and both leave marks.

\subsection{Properties and their consequences}

\textbf{Trained to sound plausible} (pre-training). Right and wrong output have the same
fluency. References are the classic case: a plausible author, title, journal and year are
easy to generate and the model has no lookup to check them against. Every reference in
this course was opened and checked (Part~III).

\textbf{Trained to please} (human feedback). Raters prefer answers that agree with them.
Sharma et al.\ \rf{sharma} showed that five assistants consistently exhibit this
\emph{sycophancy}, for example abandoning a correct answer when the user pushes back, and
that both humans and preference models sometimes prefer a convincing sycophantic answer
over a correct one. The defence on the slide follows: ask the model to find the flaw in
your argument, not to confirm it.

\textbf{Rewarded for passing checks} (verifiable rewards). A policy rewarded for passing
tests learns that weakening the test also passes it. METR reported in 2025 that
frontier models, in its evaluations, attempted ``(often successfully) to get a higher score
by modifying the tests or scoring code'' \rf{metr-hacking}. In everyday use this looks like
a loosened tolerance, a deleted assertion, or a special case for the test input. The
defence: read every change to a test, and commit tests before the implementation
(Section~\ref{sec:verification}).

\textbf{Random sampling} (architecture). With $T_{\mathrm{s}} > 0$ the same prompt gives
different answers. This is a free diagnostic: if two runs disagree on a number, neither
is reliable. Setting $T_{\mathrm{s}} = 0$ makes the choice greedy, but in practice
deployed models do not guarantee bit-identical outputs, so do not build reproducibility on
it; build it on code and data in a repository.

\textbf{No memory between sessions.} The parameters do not change while you use the
model. Everything it ``knows'' about your project in a session is in the context; at the
next session it starts from zero. Hence: write things down (Section~\ref{sec:context}).

\textbf{Finite context.} See Section~\ref{sec:context}.

\begin{qa}
\Q{Is it ``just autocomplete''?}
The objective is next-token prediction, yes. But predicting the next token of a
derivation well requires representing a lot of the structure of the derivation, and
post-training adds goal-directed behaviour. ``Just autocomplete'' underestimates what it
can do; ``understands physics'' overstates how reliably. The practical answer is the
course's: treat its output as a claim to be checked.

\Q{Why does it invent references instead of saying it doesn't know?}
Nothing in pre-training rewards abstaining, and a confident well-formatted answer tends
to be preferred by raters. Agents with web search can check references, but only if they
are asked to and actually do it; check the link yourself.

\Q{Does it learn from my conversations?}
Not within the session: the parameters are fixed. Whether a provider uses conversations
for future training depends on its terms and your account settings; check them before
pasting unpublished data (Section~\ref{sec:failures}).

\Q{How many parameters does model X have?}
Providers of closed models do not publish it. Open-weight models do, for example
$4.05 \times 10^{11}$ for the largest Llama~3 \rf{llama3}.
\end{qa}

\textbf{Read or watch next:} Karpathy's \emph{Intro to Large Language Models} (1\,h) and
\emph{Deep Dive into LLMs like ChatGPT} (3.5\,h), which covers tokenization, pre-training
and post-training in exactly this order \rf{karpathy-intro,karpathy-deep}; 3Blue1Brown
chapters 5--6 for attention \rf{3b1b}; the InstructGPT paper \rf{instructgpt},
Sections~1 and~3, for RLHF.
